\input{../header.tex}

\setlength{\parindent}{0em}

\begin{document}

\subsection{Implementation}
For the implementation of the NLL of the galah4-dataset, the same CNN network is used as in the solution from the last exercise.

The CNN increases the amount of feature maps from 1 \rightarrow 10 \rightarrow 20 \rightarrow 40 in order to achieve a deeper network, which is able to learn more complex structures in the spectrum.
After that, the features maps are compressed from a size of 40 \rightarrow 10 \rightarrow 12 \rightarrow 10, resulting in a more compact representation of the network.
In the end, the \textit{Linear}-layers are converting the features into the six \textit{Output}-neurons, representing the three labels ($T_\text{eff}, \log(g)$ and $\text{fe\_h}$) and their respective uncertainty.

For the dataset and training I used the following parameters \textit{num\_epochs} = 100 (which were not used due to early stopping), \textit{learning\_rate} = $2e-4$ and \textit{batch\_size} = 32.

\subsection{Results}
After 34 epochs, the early stopping triggered, resulting in losses of:
\begin{align*}
  \text{best\_train\_loss} &= -2.1988 (\text{in epoch 33/34})\\
  \text{best\_val\_loss} &= -2.3931 (\text{in epoch 24/34})\\
  \text{test\_loss} &= -2.2870
\end{align*}
The training loss is constantly decreasing, while the validation loss is globally decreasing, but locally oscillating. This behavior is normal and may be due to stochastic uncertainties.
Since the NLL is using a loss function of
\begin{equation}
  \nonumber
  \frac{1}{2} \left(\frac{y_i - \mu_i}{\sigma_i}\right)^2 + \log{\sigma_i},
\end{equation}
negative loss values are possible if $\sigma < 1$ due to the logarithmic term of $\log{\sigma}$ part.
Therefore, the loss values are as expected and show a good training on a CNN network with the use of early stopping.

To validate the predicted uncertainties, a pull distribution is performed. For the different labels the following mean values $\mu$ and standard deviations $\sigma$ were calculated:
\begin{align*}
  T_\text{eff} &: \mu = -0.103 &  \sigma = 0.686\\
  \log(g) &: \mu = -0.208 & \sigma = 0.822\\
  \text{fe\_h} &: \mu = -0.466 & \sigma = 0.497
\end{align*}
For very good predicted uncertainties, we expect values of a normal distribution with $\mu = 0$ and $\sigma = 1$.
All three labels have slightly shifted mean values towards the negative area. This means, that the network on average underpredicts the labels.
For the mean value the uncertainty prediction is best for the label of $T_\text{eff}$, then $\log(g)$ and lowest for the $\text{fe\_h}$.
The standard deviation on the other hand is too low in all three cases with $\log(g)$ being the closest to $\sigma=1$, followed by $T_\text{eff}$ and $\text{fe\_h}$.
Therefore, the model overestimated the uncertainty for all labels, which could be decreased by a better network/training.

Even though in some applications (e.g the safety regulations for the engine lifetime from last semester) a too conservative uncertainty is better than one being too small,
the network should aim to reach the expected distribution parameters $\mu, \sigma$.

\subsection{Problems and Improvements}
The biggest problem in this training is the instability of the standard deviation prediction $\sigma$. Due to the logarithmic part, the loss is especially sensitive to the standard deviation $\sigma$.
Training and optimizing hyperparameters become way harder, since slight changes in these parameters can have huge impacts to the loss function.
While the loss values of the validation and test seem quite good, the distribution parameters $\mu$ and $\sigma$ can still be improved for all three labels.
Some possibilities for improvement could be an adaptive \textit{learning\_rate}, which automatically changes the \textit{learning\_rate} if it isn't decreasing in a given \textit{patience}.
Another possibility lies in a grid search to optimize the hyperparameters. But since the hyperparameters are extremely sensitive in this case, this would require way more computational resources.
Since the used CNN achieved quite good results in the sole predictions of the three labels, I don't think, that a change in the network structure would improve the distribution parameters.
\end{document}
